% paper/figures_src/03_batching.tikz.tex
%
% Figure 3 — H1 tool-call batching case study (sphinx-doc__sphinx-7757,
% Codex, seed 3, both arms PASS). Structural mockup of one LLM step.
%
% Top panel (baseline arm): one assistant turn followed by three parallel
% exec_command tool calls — agent gathers context with three independent
% lookups.
%
% Bottom panel (code-only arm): same assistant turn followed by a single
% execute_code call bundling the same three lookups into one script.
%
% Intent-only labels (per Design decision §6 of figures_outline.md): cards
% carry the *purpose* of each lookup, not the literal command. File names
% appear as small grey subtitles for grounding only. No regex / Python
% source.
%
% Wrapped inside a \begin{figure}...\end{figure} float in
% sections/06_discussion.tex. main.tex must load \usepackage{tikz} +
% \usetikzlibrary{positioning,fit,backgrounds,arrows.meta}.

{% local scope for the inline-formatting commands
\providecommand{\toolhdr}[1]{{\footnotesize\sffamily\bfseries\color{black!75}#1}}
\providecommand{\intent}[1]{{\footnotesize #1}}
\providecommand{\subtitle}[1]{{\scriptsize\ttfamily\color{black!55}#1}}

\begin{tikzpicture}[
  font=\small,
  every node/.style={inner sep=2pt},
  bubble/.style={
    draw=black!50, line width=0.4pt,
    rounded corners=2pt,
    fill=black!4,
    text width=0.78\columnwidth,
    align=left,
    inner sep=4pt,
    font=\footnotesize\itshape,
  },
  toolcard/.style={
    draw=black!70, line width=0.5pt,
    rounded corners=3pt,
    text width=0.62\columnwidth,
    align=left,
    inner sep=4pt,
  },
  bundlecard/.style={
    draw=black!70, line width=0.7pt,
    rounded corners=3pt,
    text width=0.74\columnwidth,
    align=left,
    inner sep=5pt,
    fill=black!2,
  },
  armtag/.style={font=\footnotesize\sffamily\bfseries, text=black!80},
  arrow/.style={->, >=stealth, line width=0.4pt, draw=black!55},
  node distance=3pt and 0pt,
]

% =========================================================================
% TOP PANEL — baseline arm (3 parallel exec_command calls)
% =========================================================================
\node[armtag] (baseTag) at (0, 0)
  {Baseline arm \,---\, 1 LLM step \, $\rightarrow$ \, 3 tool calls};

\node[bubble, below=4pt of baseTag, anchor=north] (baseBubble)
  {Assistant: I need to see the definition, an existing test, and how
   it's called elsewhere.};

\node[toolcard, below=6pt of baseBubble.south west, anchor=north west]
  (b1) {%
    \toolhdr{exec\_command}\\
    \intent{read definition site}\\
    \subtitle{sphinx/domains/python.py}%
  };
\node[toolcard, below=3pt of b1.south west, anchor=north west]
  (b2) {%
    \toolhdr{exec\_command}\\
    \intent{read existing test}\\
    \subtitle{tests/test\_util\_inspect.py}%
  };
\node[toolcard, below=3pt of b2.south west, anchor=north west]
  (b3) {%
    \toolhdr{exec\_command}\\
    \intent{search call sites}\\
    \subtitle{tests/}%
  };

% Faint left "bracket" to suggest the three are parallel
\draw[arrow] (baseBubble.south) -- ++(0, -0.18);
\draw[black!40, line width=0.4pt]
  ([xshift=-6pt]b1.north west) -- ([xshift=-6pt]b3.south west);
\draw[black!40, line width=0.4pt]
  ([xshift=-6pt]b1.north west) -- ([xshift=-3pt]b1.north west);
\draw[black!40, line width=0.4pt]
  ([xshift=-6pt]b3.south west) -- ([xshift=-3pt]b3.south west);

% Divider between panels
\draw[black!25, line width=0.4pt, dashed]
  ([yshift=-10pt]b3.south west) -- ++(\columnwidth, 0);

% =========================================================================
% BOTTOM PANEL — code-only arm (1 execute_code call bundling all three)
% =========================================================================
\node[armtag, below=14pt of b3.south west, anchor=north west] (coTag)
  {Code-only arm \,---\, 1 LLM step \, $\rightarrow$ \, 1 tool call};

\node[bubble, below=4pt of coTag.south west, anchor=north west] (coBubble)
  {Assistant: I need to see the definition, an existing test, and how
   it's called elsewhere.};

\node[bundlecard, below=6pt of coBubble.south, anchor=north]
  (co1) {%
    \toolhdr{execute\_code} \, \subtitle{one script}\\[2pt]
    \intent{$\bullet$\, read definition site} \hfill \subtitle{sphinx/domains/python.py}\\
    \intent{$\bullet$\, read existing test} \hfill \subtitle{tests/test\_util\_inspect.py}\\
    \intent{$\bullet$\, search call sites} \hfill \subtitle{tests/}%
  };

\draw[arrow] (coBubble.south) -- (co1.north);

\end{tikzpicture}
}% end local scope
